%% ============================================================
%% Introduction to Cognitive Psychology — Week 4 Study Notes
%% ============================================================
\documentclass[11pt,a4paper]{article}

\usepackage[a4paper, left=1.8cm, right=1.8cm, top=1.3cm, bottom=1.8cm]{geometry}
\usepackage{booktabs}
\usepackage{tabularx}
\usepackage{array}
\usepackage{enumitem}
\usepackage{titlesec}
\usepackage{fancyhdr}
\usepackage{fontenc}
\usepackage{inputenc}
\usepackage{microtype}
\usepackage{parskip}
\usepackage{hyperref}
\usepackage{amsmath}
\usepackage{amssymb}

\hypersetup{hidelinks, colorlinks=false}

%% ── Table helpers ─────────────────────────────────────────
\newcolumntype{L}[1]{>{\raggedright\arraybackslash}p{#1}}

%% ── Section headings ──────────────────────────────────────
\titleformat{\section}{\large\bfseries}{}{0em}{}[\titlerule]
\titleformat{\subsection}{\normalsize\bfseries}{}{0em}{}
\titleformat{\subsubsection}{\normalsize\bfseries}{}{0em}{}

\titlespacing{\section}{0pt}{12pt}{4pt}
\titlespacing{\subsection}{0pt}{8pt}{3pt}
\titlespacing{\subsubsection}{0pt}{6pt}{2pt}

%% ── Page style ────────────────────────────────────────────
\pagestyle{fancy}
\fancyhf{}
\renewcommand{\headrulewidth}{0pt}
\fancyfoot[C]{\small Cognitive Psychology --- Week 4 \textbar\ Page \thepage}

%% ── Misc helpers ──────────────────────────────────────────
\newcommand{\bb}[1]{\textbf{#1}}
\setlength{\parskip}{4pt}

%% ═══════════════════════════════════════════════════════════
\begin{document}
\setlength{\arrayrulewidth}{0.4pt}

\begin{center}
\bfseries\LARGE Introduction to Cognitive Psychology\par
\vspace{0.2cm}
\large Assignment 4 Detailed Solutions
\end{center}
\vspace{0.5cm}

%% ════════════════════════════════════════════════════════════
\section*{ASSIGNMENT 4 --- Full Solutions with Explanations}

\subsubsection*{Question 1}
\textit{Which of the following is NOT a component of Baddeley's working memory model?}
\medskip

\begin{tabular}{L{0.3cm} L{14.9cm}}
& A.\ The phonological loop \\
& \textbf{B.\ The icon \checkmark} \\
& C.\ The central executive \\
& D.\ The visuospatial sketchpad \\
\end{tabular}

\vspace{0.3cm}
\noindent\textbf{Ans:} B — The icon\\[0.2cm]
\textbf{Why this answer:} \\
\textbf{Baddeley's working memory model} consists of three original components: the \textbf{phonological loop} (verbal/acoustic information), the \textbf{visuospatial sketchpad} (spatial and visual information), and the \textbf{central executive} (attentional control system). A fourth component, the \textbf{episodic buffer}, was added later (2000). The \textbf{icon} is not part of this model — it refers to the brief visual trace stored in \textbf{iconic (sensory) memory}, which is a separate memory system that precedes working memory.
\vspace{0.4cm}

\subsubsection*{Question 2}
\textit{Sternberg's classic work on searching for information from short-term memory indicated that the search process is:}
\medskip

\begin{tabular}{L{0.3cm} L{14.9cm}}
& A.\ Serial \\
& B.\ Self-terminating \\
& C.\ Parallel \\
& \textbf{D.\ Both serial and exhaustive \checkmark} \\
\end{tabular}

\vspace{0.3cm}
\noindent\textbf{Ans:} D — Both serial and exhaustive\\[0.2cm]
\textbf{Why this answer:} \\
In his classic memory-scanning experiments, \textbf{Saul Sternberg} (1966) found that reaction time increased linearly with the number of items in the memory set, regardless of whether the item was present or absent. This pattern of results supports a \textbf{serial exhaustive search} model — meaning the search proceeds \textbf{item by item} (serial) and continues through the \textbf{entire set} even after a match has been found (exhaustive), rather than stopping immediately upon finding the target (self-terminating).
\vspace{0.4cm}

\subsubsection*{Question 3}
\textit{In the absence of rehearsal, short-term memory tends to:}
\medskip

\begin{tabular}{L{0.3cm} L{14.9cm}}
& \textbf{A.\ Last about 20 seconds \checkmark} \\
& B.\ Last about 8 seconds \\
& C.\ Decay slowly over 24 hours \\
& D.\ Decay slowly over a week \\
\end{tabular}

\vspace{0.3cm}
\noindent\textbf{Ans:} A — Last about 20 seconds\\[0.2cm]
\textbf{Why this answer:} \\
Classic research by \textbf{Peterson and Peterson (1959)} demonstrated that without active rehearsal, information in short-term memory decays rapidly within approximately \textbf{15–30 seconds} (commonly cited as ~18–20 seconds). Their paradigm used consonant trigrams followed by a distractor task (counting backward) to prevent rehearsal, showing that recall dropped dramatically within this brief window — far shorter than what would be expected from long-term memory decay.
\vspace{0.4cm}
\newpage

\subsubsection*{Question 4}
\textit{The suffix effect relates to which type of memory?}
\medskip

\begin{tabular}{L{0.3cm} L{14.9cm}}
& A.\ Iconic \\
& \textbf{B.\ Echoic \checkmark} \\
& C.\ Short term \\
& D.\ Working \\
\end{tabular}

\vspace{0.3cm}
\noindent\textbf{Ans:} B — Echoic\\[0.2cm]
\textbf{Why this answer:} \\
The \textbf{suffix effect} occurs when an irrelevant auditory stimulus (a "suffix") presented at the end of a list disrupts recall of the final items. This specifically demonstrates the brief persistence of \textbf{echoic memory} — the auditory sensory store that holds sound-based information for approximately 2–4 seconds. The suffix overwrites or interferes with the echoic trace of the last list items, reducing the recency advantage normally observed in serial recall.
\vspace{0.4cm}

\subsubsection*{Question 5}
\textit{Unattended information is stored briefly in:}
\medskip

\begin{tabular}{L{0.3cm} L{14.9cm}}
& \textbf{A.\ Sensory memory \checkmark} \\
& B.\ Short term memory \\
& C.\ Long term memory \\
& D.\ Working memory \\
\end{tabular}

\vspace{0.3cm}
\noindent\textbf{Ans:} A — Sensory memory\\[0.2cm]
\textbf{Why this answer:} \\
\textbf{Sensory memory} is the initial, very brief storage system that holds all incoming sensory information — both attended and unattended — for a fraction of a second (iconic/visual) to a few seconds (echoic/auditory). Only information that receives \textbf{attentional selection} is transferred to short-term or working memory for further processing. Unattended information remains briefly in sensory memory before it rapidly fades and is lost.
\vspace{0.4cm}

\subsubsection*{Question 6}
\textit{Retrieval involves:}
\medskip

\begin{tabular}{L{0.3cm} L{14.9cm}}
& A.\ The activation of the senses \\
& B.\ The translation of information into a form that can be stored \\
& C.\ The storage of information over time \\
& \textbf{D.\ The calling to mind of previously stored information \checkmark} \\
\end{tabular}

\vspace{0.3cm}
\noindent\textbf{Ans:} D — The calling to mind of previously stored information\\[0.2cm]
\textbf{Why this answer:} \\
\textbf{Retrieval} is the process of accessing and bringing previously stored information back into conscious awareness from long-term memory. It is the third stage of the memory process, following \textbf{encoding} (translating information into a storable form — option B) and \textbf{storage} (maintaining information over time — option C). Retrieval can occur through various mechanisms including \textbf{recall} (generating information from memory), \textbf{recognition} (identifying previously encountered items), and \textbf{relearning}.
\vspace{0.4cm}
\newpage

\subsubsection*{Question 7}
\textit{The primacy and recency effects in memory:}
\medskip

\begin{tabular}{L{0.3cm} L{14.9cm}}
& A.\ Are thought to be due to the action of short-term memory \\
& B.\ Are thought to be due to the action of long-term memory \\
& C.\ Are thought to be due to the action of sensory memory \\
& \textbf{D.\ Can be independently manipulated, indicating at least two types of memory at work \checkmark} \\
\end{tabular}

\vspace{0.3cm}
\noindent\textbf{Ans:} D — Can be independently manipulated, indicating at least two types of memory at work\\[0.2cm]
\textbf{Why this answer:} \\
The \textbf{serial position curve} shows two distinct effects: the \textbf{primacy effect} (better recall of early items, attributed to \textbf{long-term memory} due to greater rehearsal) and the \textbf{recency effect} (better recall of final items, attributed to \textbf{short-term memory} since they are still in the active buffer). Critically, these effects can be \textbf{independently manipulated} — for example, a delay with a distractor task eliminates the recency effect while leaving the primacy effect intact. This double dissociation provides strong evidence for the existence of at least \textbf{two separate memory systems}.
\vspace{0.4cm}

\subsubsection*{Question 8}
\textit{Information in short-term memory is assumed to be coded primarily by:}
\medskip

\begin{tabular}{L{0.3cm} L{14.9cm}}
& \textbf{A.\ Sound \checkmark} \\
& B.\ Meaning \\
& C.\ Visual appearance \\
& D.\ Both sound and meaning \\
\end{tabular}

\vspace{0.3cm}
\noindent\textbf{Ans:} A — Sound\\[0.2cm]
\textbf{Why this answer:} \\
Research by \textbf{Conrad (1964)} demonstrated that short-term memory relies primarily on \textbf{acoustic (phonological) coding}. He found that participants made more errors confusing letters that sounded similar (e.g., B, C, D, G, P, T, V) than letters that looked similar, even when the letters were presented visually. This contrasts with \textbf{long-term memory}, which relies more heavily on \textbf{semantic (meaning-based) coding}. This distinction in coding formats is one of the key differences between STM and LTM.
\vspace{0.4cm}

\subsubsection*{Question 9}
\textit{Repeating a phone number to yourself to hold it in memory while you dial it would use which component of working memory?}
\medskip

\begin{tabular}{L{0.3cm} L{14.9cm}}
& A.\ The visuospatial sketchpad \\
& \textbf{B.\ The phonological loop \checkmark} \\
& C.\ The episodic buffer \\
& D.\ Both the visuospatial sketchpad and the phonological loop \\
\end{tabular}

\vspace{0.3cm}
\noindent\textbf{Ans:} B — The phonological loop\\[0.2cm]
\textbf{Why this answer:} \\
The \textbf{phonological loop} is the component of Baddeley's working memory model responsible for maintaining \textbf{verbal and acoustic information} through a process of \textbf{subvocal rehearsal} (inner speech). It consists of two sub-components: the \textbf{phonological store} (which holds sound-based representations for ~2 seconds) and the \textbf{articulatory rehearsal process} (which refreshes these representations through repetition). Silently repeating a phone number is a classic example of the articulatory rehearsal process keeping information active in the phonological loop.
\vspace{0.4cm}
\newpage
\subsubsection*{Question 10}
\textit{Disrupting the process of long-term potentiation leads to:}
\medskip

\begin{tabular}{L{0.3cm} L{14.9cm}}
& \textbf{A.\ Disruption of learning and remembering \checkmark} \\
& B.\ Recovery from anterograde amnesia \\
& C.\ Recovery from retrograde amnesia \\
& D.\ Language deficits \\
\end{tabular}

\vspace{0.3cm}
\noindent\textbf{Ans:} A — Disruption of learning and remembering\\[0.2cm]
\textbf{Why this answer:} \\
\textbf{Long-term potentiation (LTP)} is a neurophysiological process in which repeated stimulation of synapses leads to a \textbf{lasting strengthening of synaptic connections}. LTP is widely regarded as one of the primary neural mechanisms underlying \textbf{learning and memory formation}, particularly in the \textbf{hippocampus}. Disrupting LTP — through pharmacological agents, genetic manipulation, or lesions — interferes with the ability to form new memories and learn new information, confirming its critical role in memory consolidation.
\vspace{0.4cm}

\medskip
\subsection*{Assignment Quick Answer Summary}

\begin{center}
\renewcommand{\arraystretch}{1.4}
\begin{tabular}{|L{0.8cm}|L{6.2cm}|L{8.2cm}|}
\hline
\bb{Q\#} & \bb{Topic} & \bb{Correct Answer} \\
\hline
1 & Baddeley's Working Memory Model & B — The icon (not a component) \\
2 & Sternberg's Memory Search & D — Serial and exhaustive \\
3 & STM Duration Without Rehearsal & A — About 20 seconds \\
4 & Suffix Effect & B — Echoic memory \\
5 & Storage of Unattended Information & A — Sensory memory \\
6 & Definition of Retrieval & D — Calling to mind stored info \\
7 & Primacy \& Recency Effects & D — Independently manipulated (two systems) \\
8 & STM Coding Format & A — Sound (acoustic coding) \\
9 & Phonological Loop & B — Phonological loop \\
10 & Long-Term Potentiation & A — Disruption of learning and remembering \\
\hline
\end{tabular}
\end{center}

\medskip
\end{document}
